\chapter{大语言模型生成内容安全增强系统设计与开发}\label{chap:system}

本章围绕鲁棒安全编辑和多类别动态路由的应用需求，设计模型登记、编辑执行、模块管理、交互推理和结果追溯一体化流程。系统通过统一配置与评估接口连接算法和应用，支持安全干预过程的可管理、可比较和可复核。当前为需求与设计框架，代码、界面、联调及测试\textbf{待完成}。

\section{需求分析}
\subsection{用户需求}

系统面向算法研究与演示使用者，支持登记本地模型、配置鲁棒安全编辑和多类别路由任务、比较编辑前后及实际LoRA微调后的输出，并查看路由依据与安全效用指标。界面明确区分任务提交、运行、完成及失败状态。

\subsection{功能需求}

拟提供模型与数据版本登记、编辑配置、离线任务、检查点与模块管理、交互推理、批量评估和结果导出。第二章流程支持子空间缓存及LoRA代理库；第三章支持分类头、等级阈值和类别模块绑定。汇总表须可追溯至原始输出和指标配置。

\subsection{非功能需求}

优先保证结果可复核、任务可恢复和资源可控。拟采用单GPU串行任务队列，避免编辑、微调与推理同时争用显存；路由状态与请求绑定。配置错误、显存不足和版本不兼容应返回明确原因，不静默切换模型。日志按需脱敏保存。时长、显存与吞吐量需实测，不预写性能承诺。范围限定本地研究原型，不将集群、多租户及复杂权限平台列为必做工作。

\section{系统概要介绍}

拟采用Python和PyTorch组织模型运行，通过Transformers统一模型与分词器的加载接口\cite{wolf2020transformers}。低秩适配由PEFT层管理，编辑任务及基准评估以EasyEdit为实验基础设施\cite{hu2022lora,wang2024easyedit}。实际依赖版本在实现后锁定。交互层拟采用轻量界面，实验记录使用SQLite与JSON，模型、子空间和类别模块使用版本化文件存储。总体架构如图\ref{fig:system-architecture}所示。

\begin{figure}[htbp]
 \centering
 \fbox{\begin{minipage}{0.88\textwidth}\centering
 交互层：模型配置、编辑任务、推理对比、评估报告\\[3pt]
 $\downarrow$\\[3pt]
 服务层：任务队列、配置校验、版本绑定、结果汇总\\[3pt]
 $\downarrow$\\[3pt]
 算法层：鲁棒安全编辑 / 类别分级路由 / 基准方法评估\\[3pt]
 $\downarrow$\\[3pt]
 存储层：基座、子空间、LoRA代理库、模块及逐样本结果
 \end{minipage}}
 \caption{安全增强系统拟采用的总体架构（待开发）}
 \label{fig:system-architecture}
\end{figure}

两项方法默认作为独立模式：第二章加载编辑检查点，可继续真实LoRA微调；第三章在其绑定基座上加载类别模块与路由器。不能把基于$\theta_0$的模块直接叠加到第二章$\theta_e$并假定效果不变。若增加组合演示，需重新校准或训练并独立评估；跨章叠加不作为基本完成要求。

\section{功能模块设计}
\subsection{基础模型管理模块}

登记模型及分词器版本、聊天模板、精度和加载配置，用文件及配置摘要识别基座，检查编辑层、模块参数名与维度。所有产物关联到确定基座版本，加载与卸载更新资源占用状态。

\subsection{鲁棒安全编辑模块}

封装局部编辑位置确定、参数冻结、$Q_s/Q_u$构造、LoRA代理加载、式\eqref{eq:full-objective}优化及检查点输出。缓存记录数据摘要、参数坐标、秩和归一化规则；代理库记录真实LoRA的数据、步数、目标层及种子。界面区分固定代理叠加与从编辑模型实际微调，缺少真实日志时保持率显示“待完成”。

\subsection{多类别安全编辑模块}

维护类别模块的创建、保存、加载和禁用，绑定类别、秩、缩放、编辑层及基座标识。新增模块后评估目标类、其他类别及无害集；版本更新保留历史结果，便于追踪干扰矩阵变化。

\subsection{风险感知动态路由模块}

加载匹配基座的分类头及锁定阈值，按式\eqref{eq:ch3-risk-gate}和式\eqref{eq:ch3-route}执行门控。展示等级、类别、模块与回退标记；请求结束后清除路由状态。原型优先串行推理，避免请求间状态污染。

\subsection{安全响应生成模块}

统一原始模型、DINM及两章方法的生成参数，支持同提示对比。长度、采样、种子和模板写入记录。回退触发时明确标记，保留算法本身的独立评估记录，不能用界面回退效果替代模型测试。

\subsection{安全性能评估模块}

计算ASR、SRR、真实微调后保持率、效用、ORR、类别干扰和路由指标，导出逐样本与汇总数据。记录判别器版本、人工复核及分歧规则；未生成的结果显示“待测”，不使用随机演示数值冒充实验。

\section{系统开发}

开发计划依次覆盖模型与数据登记、编辑任务执行、真实LoRA评估、类别模块与路由集成、结果导出及交互界面。提交编辑任务时选择已登记的模型和数据版本，服务先校验配置，再创建任务记录并排队执行。任务成功后登记检查点或模块及其来源。失败时保留配置与失败原因，支持用同一配置重新提交。一次对比评估使用同一测试清单和生成设置，汇总指标关联到逐样本记录。表\ref{tab:system-records}列出主要记录字段。

\begin{table}[htbp]
 \centering
 \caption{系统核心记录设计（实现待完成）}
 \label{tab:system-records}
 \begin{tabular}{p{0.2\textwidth}p{0.65\textwidth}}
 \toprule
 记录 & 关键字段\\ \midrule
 模型与数据 & 标识、版本摘要、模板、数据用途与划分清单\\
 编辑任务 & 方法、基座、配置、种子、状态、产物及失败原因\\
 编辑产物 & 检查点、子空间或模块、坐标/维度、来源任务\\
 推理记录 & 输入、输出、参数、等级、模块及回退标记\\
 评估记录 & 测试清单、判别器、逐样本标签、指标及复核\\
 \bottomrule
 \end{tabular}
\end{table}

实际依赖版本、接口实现、一次完整运行记录与界面截图待补充。

\section{系统测试}

系统测试按表\ref{tab:system-tests}验证功能与集成。对版本不匹配用例，预期任务进入失败状态并保留提示，相关组件不进入运行阶段。对路由隔离用例，连续提交激活请求和旁路请求，检查第二次请求的调用记录及权重是否符合旁路规则。对结果复核用例，由导出的逐样本数据独立汇总指标，与系统报告逐项比对。测试结果关联硬件、配置和日志，当前待执行。
\begin{table}[htbp]
 \centering
 \caption{系统测试计划及状态}
 \label{tab:system-tests}
 \begin{tabular}{p{0.17\textwidth}p{0.53\textwidth}p{0.14\textwidth}}
 \toprule
 维度 & 验证内容 & 状态\\ \midrule
 模型与模块 & 正确加载、不兼容阻断及卸载资源释放 & 待完成\\
 鲁棒编辑 & 子空间/代理追溯，真实LoRA评估独立执行 & 待完成\\
 类别路由 & 三级门控、类别切换及请求后状态清除 & 待完成\\
 端到端流程 & 配置、编辑、推理、评估与导出一致 & 待完成\\
 异常处理 & 显存不足、中断、模块缺失及配置错误 & 待完成\\
 性能 & 时长、峰值显存、首token延迟、吞吐量 & 待测\\
 结果复核 & 逐样本记录可复算指标，回退指标单列 & 待完成\\
 \bottomrule
 \end{tabular}
\end{table}

实施后补充失败项、修复及复测记录。性能测试固定模型、精度、输入/输出长度和硬件，说明预热和重复次数；若只支持串行请求，应据实说明边界。

\section{本章小结}

本章形成需求、架构、六类功能模块及测试方案，为鲁棒编辑和类别路由提供统一任务与评估入口。范围为单GPU本地原型，优先完成可复现闭环。开发、截图、性能与集成结论均\textbf{待完成}，完成后再撰写工程结果。
